% -*- coding: utf-8 -*-
\documentclass[12pt,a4paper]{report}
\usepackage[utf8]{inputenc}
\usepackage[T1]{fontenc}
\usepackage[french]{babel}
\usepackage{microtype}
\usepackage{geometry}
\geometry{left=3cm,right=2.5cm,top=2.5cm,bottom=2.5cm}
\usepackage{graphicx}
\usepackage{caption}
\usepackage{subcaption}
\usepackage{float}
\usepackage{booktabs}
\usepackage{longtable}
\usepackage{hyperref}
\hypersetup{colorlinks=true,linkcolor=blue,citecolor=blue,urlcolor=blue}
\usepackage{fancyhdr}
\usepackage{tocloft}
\usepackage{listings}
\usepackage{xcolor}
\usepackage{enumitem}

% Chemin par défaut pour les images
\graphicspath{{figures/}}

% Header / Footer
\pagestyle{fancy}
\fancyhf{}
\rhead{Projet Chat-AI}
\lhead{Rapport Technique}
\cfoot{\thepage}

% Styles pour le code
\definecolor{codegray}{rgb}{0.95,0.95,0.95}
\lstset{
  backgroundcolor=\color{codegray},
  basicstyle=\ttfamily\footnotesize,
  breaklines=true,
  frame=single,
  numbers=left,
  numberstyle=\tiny,
  captionpos=b
}

% Metadonnées du document
\title{\Huge{Développement d'une application de type Chat-AI \\ (Frontend + Backend)}}
\author{Équipe de développement \\ Réalisé par : Votre Nom \\ Encadrant : Nom de l'encadrant}
\date{\today}

\begin{document}

% Page de garde
\begin{titlepage}
  \centering
  \vspace*{2cm}
  {\Huge \textbf{Développement d'une application Chat-AI}}\\[1cm]
  {\Large Frontend (React/Vite/TypeScript) \& Backend (Java Spring Boot)}\\[2cm]
  \includegraphics[width=0.45\textwidth]{cover-logo.png}\\[1.5cm]
  {\Large \textbf{Rapport Technique et Académique}}\\[2cm]
  {\large Réalisé par : \textbf{Votre Nom}}\\
  {\large Équipe : \textbf{Nom de l'équipe / Auteur}}\\[1cm]
  {\large Encadrant : \textbf{Nom de l'encadrant}}\\[2cm]
  {\large Date : \today}\\
  \vfill
  {\small Université / École — Département / Master / Licence}
\end{titlepage}

% Résumé en français
\begin{abstract}
  Ce rapport présente la conception, l'implémentation et le déploiement d'une application de type Chat-AI comprenant une interface web (frontend) développée en React + TypeScript et un serveur backend Java Spring Boot. Le document couvre l'architecture logicielle, les choix technologiques, les détails d'implémentation clés pour les composants critiques, la sécurité, le déploiement Docker / docker-compose, ainsi que des suggestions pour des améliorations futures. Des emplacements pour captures d'écran et diagrammes sont prévus.
\end{abstract}

\clearpage
\tableofcontents
\listoffigures
\listoftables
\clearpage

% Remerciements (facultatif)
\chapter*{Remerciements}
\addcontentsline{toc}{chapter}{Remerciements}
Nous remercions... (Compléter selon le contexte).

\clearpage

% Introduction
\chapter{Introduction}
\section{Contexte}
Présentez le contexte général du projet : assistants conversationnels, intégration d'IA, motivations (ex : aide à la productivité, interfaces vocales, accessibilité, etc.).

\section{Objectifs}
\begin{itemize}
  \item Concevoir une application web de chat assistée par IA.
  \item Fournir une architecture scalable et sécurisée.
  \item Intégrer TTS (Text-To-Speech) et gestion de fichiers (PDF).
  \item Déployer via Docker / docker-compose.
\end{itemize}

\section{Organisation du rapport}
Brève description des chapitres.

% Chapitre architecture
\chapter{Analyse et Architecture du Système}
\section{Vue d'ensemble}
Décrire l'architecture globale : client web (React) $\\leftrightarrow$ API REST + WebSocket (Spring Boot) $\\leftrightarrow$ Services externes (API Gemini, TTS, Redis).

\begin{figure}[H]
  \centering
  \includegraphics[width=0.85\textwidth]{architecture.png}
  \caption{Architecture générale du système (client, serveur, services externes).}
  \label{fig:architecture}
\end{figure}

\subsection{Composants principaux}
\begin{itemize}
  \item Frontend : \texttt{ask-ai-web-app} (Vite, React, TypeScript)
  \item Backend : \texttt{ai-server} (Spring Boot, Maven)
  \item Persistances : Redis, base de données (si applicable)
  \item Déploiement : Docker / docker-compose
\end{itemize}

\subsection{Diagrammes recommandés}
\begin{itemize}
  \item Diagramme d'architecture (figure \ref{fig:architecture}).
  \item Diagramme ER pour les entités utilisateur / messages (\texttt{figures/er_diagram.png}).
  \item Diagramme de séquence pour l'envoi d'un message et réponse du modèle IA (\texttt{figures/sequence_message.png}).
\end{itemize}

% Chapitre Backend
\chapter{Backend --- Spring Boot}
\section{Présentation du module}
Chemin du projet : \texttt{backend/ai-server/}\\
Fichier de configuration principal : \texttt{src/main/resources/application.yml}

\section{Technologies et dépendances}
\begin{itemize}
  \item Java 17+ (ou la version utilisée)
  \item Spring Boot (web, security, data, websocket)
  \item Maven (wrapper fourni)
  \item Redis pour sessions / cache
  \item Bibliothèques pour TTS (client MaryTTS) et intégration API IA
\end{itemize}

\section{Structure du code}
\begin{itemize}
  \item \texttt{com.bagga.aiserver.AiServerApplication} : point d'entrée.
  \item \texttt{controller/} : endpoints REST et WebSocket (ex: \texttt{GeminiController}, \texttt{PdfController}, \texttt{TextToSpeechController}).
  \item \texttt{services/} : logique métier (\texttt{ChatMessageService}, \texttt{TextToSpeechService}, \texttt{PdfExtractorService}).
  \item \texttt{repositories/} : interfaces Spring Data (\texttt{UserRepository}, \texttt{MessageRepository}).
  \item \texttt{config/} : configuration de sécurité et JWT (\texttt{SecurityConfiguration}, \texttt{JwtService}, \texttt{JwtAuthenticationFilter}), WebSocket (\texttt{WebSocketConfig}) et Redis (\texttt{RedisConfig}).
\end{itemize}

\subsection{Endpoints clés}
\begin{itemize}
  \item \texttt{POST /api/auth/login} --- authentification
  \item \texttt{POST /api/users} --- création d'utilisateur
  \item \texttt{POST /api/gemini/ask} --- envoyer un message au modèle IA
  \item \texttt{POST /api/pdf/upload} --- extraire texte d'un PDF
  \item WebSocket: \texttt{/ws/chat} --- canal pour messages en temps réel
\end{itemize}

\subsection{Sécurité}
\begin{itemize}
  \item Authentification JWT
  \item Filtres (\texttt{JwtAuthenticationFilter}) pour valider les requêtes
  \item Bonnes pratiques : stocker secrets dans variables d'environnement, rotation des clés, limiter scopes d'API externes
\end{itemize}

\subsection{Extraits de code}
Exemple : configuration WebSocket (fichier \texttt{WebSocketConfig.java}, résumé)
\begin{lstlisting}[language=Java,caption={Extrait : configuration WebSocket}]
@Configuration
@EnableWebSocketMessageBroker
public class WebSocketConfig implements WebSocketMessageBrokerConfigurer {
  @Override
  public void registerStompEndpoints(StompEndpointRegistry registry) {
    registry.addEndpoint("/ws/chat").setAllowedOrigins("*");
  }
  // ...
}
\end{lstlisting}

% Chapitre Frontend
\chapter{Frontend --- React + TypeScript}
\section{Présentation du module}
Chemin du projet : \texttt{frontend/ask-ai-web-app/}

\section{Technologies}
\begin{itemize}
  \item Vite + React
  \item TypeScript
  \item Axios (interceptor \texttt{Interceptors/axiosInstance.ts})
  \item WebSocket client (service \texttt{services/WebSocketService.ts})
  \item Gestion d'état via Context API (\texttt{context/ChatContext.tsx})
\end{itemize}

\section{Structure du projet}
\begin{itemize}
  \item \texttt{src/pages/} : pages principales (\texttt{ChatAsistant.tsx}, \texttt{LoginPage.tsx})
  \item \texttt{src/components/} : composants réutilisables (\texttt{Chat.tsx}, \texttt{ChatAsistantComponent.tsx}, \texttt{Login.tsx})
  \item \texttt{src/apis/} : appels HTTP et DTOs
  \item \texttt{src/services/WebSocketService.ts} : connexion en temps réel
\end{itemize}

\section{Flux de données et logique}
\begin{itemize}
  \item Authentification : stockage du token JWT (localStorage ou cookie sécurisé)
  \item Connexion WebSocket établie après authentification
  \item Envoi d'un message : UI -> WebSocket ou POST -> Backend -> modèle IA -> réponse -> UI
\end{itemize}

\section{Extraits de code}
Exemple d'appel API (Axios interceptor)
\begin{lstlisting}[language=JavaScript,caption={Extrait : Interceptor Axios}]
import axios from 'axios';
const instance = axios.create({ baseURL: '/api' });
instance.interceptors.request.use((config) => {
  const token = localStorage.getItem('token');
  if (token) config.headers.Authorization = `Bearer ${token}`;
  return config;
});
export default instance;
\end{lstlisting}

\section{Captures d'écran recommandées}
\begin{itemize}
  \item \texttt{figures/ui_login.png} --- page de connexion.
  \item \texttt{figures/ui_chat.png} --- interface de chat principal.
  \item \texttt{figures/ui_test_audio.png} --- page TestAudioPage (ex. voix/TTS).
  \item \texttt{figures/ui_admin_dashboard.png} --- si existant, page d'administration.
\end{itemize}

% Chapitre Intégration & Déploiement
\chapter{Intégration et Déploiement}
\section{Fichier \texttt{docker-compose.yml}}
Inclure un extrait et expliquer les services (app-backend, app-frontend, redis, db). Exemple :
\begin{lstlisting}[language=YAML,caption={Extrait : docker-compose.yml}]
version: '3.8'
services:
  ai-server:
    build: ./ai-server
    ports:
      - "8080:8080"
    environment:
      - SPRING_PROFILES_ACTIVE=prod
  frontend:
    build: ./ask-ai-web-app
    ports:
      - "3000:3000"
  redis:
    image: redis:6
    ports:
      - "6379:6379"
\end{lstlisting}

\section{Étapes de déploiement local}
\begin{enumerate}
  \item Construire backend : \texttt{cd backend/ai-server && ./mvnw clean package}
  \item Construire images Docker et démarrer : \texttt{docker-compose up --build}
  \item Accéder à \texttt{http://localhost:3000} (frontend) et \texttt{http://localhost:8080} (API)
\end{enumerate}

\section{CI/CD (suggestion)}
Pipeline GitHub Actions / GitLab CI : build backend, run tests, build frontend, publish images, déployer.

% Chapitre Tests et Qualité
\chapter{Tests, Validation et Qualité}
\section{Tests Backend}
\begin{itemize}
  \item Tests unitaires JUnit (ex : \texttt{AiServerApplicationTests.java})
  \item Tests d'intégration pour contrôleurs et répositories
  \item Mock des services externes (Gemini API, TTS)
\end{itemize}

\section{Tests Frontend}
\begin{itemize}
  \item Tests unitaires (Jest + React Testing Library)
  \item Tests end-to-end (Cypress) pour scénarios critiques (login, envoi message, playback TTS)
\end{itemize}

\section{Analyse de qualité}
\begin{itemize}
  \item Linters (ESLint) et formatters (Prettier) pour le frontend
  \item Static analysis (SpotBugs / PMD) pour Java
\end{itemize}

% Chapitre Résultats et démonstration
\chapter{Résultats et démonstrations}
Présenter exemples d'exécutions, logs et métriques.

\begin{figure}[H]
  \centering
  \includegraphics[width=0.8\textwidth]{ui_chat.png}
  \caption{Capture d'écran: Interface Chat}
  \label{fig:ui_chat}
\end{figure}

\subsection{Scénarios testés}
\begin{itemize}
  \item Authentification et création d'utilisateur
  \item Conversation basique avec l'IA
  \item TTS de messages (génération audio)
  \item Extraction de texte depuis un PDF
\end{itemize}

% Chapitre Sécurité et confidentialité
\chapter{Sécurité et Confidentialité}
\begin{itemize}
  \item Gestion des tokens JWT et durées d'expiration
  \item Stockage des secrets (ne pas committer \texttt{application.yml} contenant secrets)
  \item Consentement et anonymisation des messages (si données utilisateurs sensibles)
  \item Recommandation : utilisation d'un vault (HashiCorp Vault/Azure Key Vault)
\end{itemize}

% Chapitre Travaux futurs
\chapter{Perspectives et améliorations}
\begin{itemize}
  \item Améliorer la scalabilité (Kubernetes)
  \item Ajouter monitoring (Prometheus + Grafana)
  \item Tester robustesse (fuzzing, tests de charge)
  \item Mode conversation multi-utilisateurs / context windows plus longues
\end{itemize}

% Références / Bibliographie
\clearpage
\chapter*{Références}
\addcontentsline{toc}{chapter}{Références}
\nocite{*}
\bibliographystyle{plain}
\bibliography{references}

% Annexes
\appendix
\chapter{Annexes}
\section{Fichiers de configuration importants}
\begin{itemize}
  \item \texttt{application.yml} (expliquer les clés)
  \item \texttt{docker-compose.yml} (extrait)
\end{itemize}

\section{Liste des fichiers clés du projet}
\begin{itemize}
  \item Backend : \texttt{src/main/java/...} (lister les packages et classes importantes)
  \item Frontend : \texttt{src/components}, \texttt{src/pages}, \texttt{src/services}
\end{itemize}

\section{Instructions rapides pour reproduire}
\begin{enumerate}
  \item Cloner le dépôt.
  \item Backend : \texttt{cd backend/ai-server && ./mvnw clean package}
  \item Frontend : \texttt{cd frontend/ask-ai-web-app && npm install && npm run build}
  \item \texttt{docker-compose up --build}
\end{enumerate}

\end{document}
